\documentclass[11pt,reqno]{amsart}

\usepackage[letterpaper,margin=1in]{geometry}
\usepackage{amsmath,amssymb,mathtools,mathrsfs}
\usepackage{booktabs,tabularx,array}
\usepackage{enumitem}
\usepackage{microtype}
\usepackage{xcolor}
\usepackage[colorlinks=true,linkcolor=blue!55!black,citecolor=blue!55!black,urlcolor=blue!55!black]{hyperref}

\newtheorem{theorem}{Theorem}[section]
\newtheorem{proposition}[theorem]{Proposition}
\newtheorem{lemma}[theorem]{Lemma}
\newtheorem{corollary}[theorem]{Corollary}
\theoremstyle{definition}
\newtheorem{definition}[theorem]{Definition}
\theoremstyle{remark}
\newtheorem{remark}[theorem]{Remark}

\newcommand{\C}{\mathbb C}
\newcommand{\Z}{\mathbb Z}
\newcommand{\cD}{\mathcal D}
\newcommand{\cK}{\mathcal K}
\newcommand{\cT}{\mathcal T}
\newcommand{\cC}{\mathcal C}
\newcommand{\cO}{\mathcal O}
\newcommand{\fD}{\mathfrak D}
\newcommand{\fK}{\mathfrak K}
\newcommand{\one}{\mathbf 1}
\newcommand{\im}{\operatorname{im}}
\newcommand{\Syz}{\operatorname{Syz}}
\newcommand{\Res}{\operatorname*{Res}}

\title[Nonsurjectivity and fibers of the Sinh--Gordon K-transform]{Nonsurjectivity and complete fibers of the Sinh--Gordon K-transform}
\author{Zhiyuan Yao}
\address{Lanzhou Center for Theoretical Physics, Key Laboratory of Theoretical Physics of Gansu Province, Key Laboratory of Quantum Theory and Applications of MoE, Gansu Provincial Research Center for Basic Disciplines of Quantum Physics, Lanzhou University, Lanzhou 730000, China}
\email{yaozy@lzu.edu.cn}
\date{\today}

\hypersetup{
  pdftitle={Nonsurjectivity and complete fibers of the Sinh-Gordon K-transform},
  pdfauthor={Zhiyuan Yao},
  pdfsubject={A pure-A4 obstruction and a coherent inverse-limit classification of fixed-parameter K-transform fibers}
}

\begin{document}

\begin{abstract}
The binary-labelled K-transform is a standard device for constructing
Sinh--Gordon form factors from entire periodic auxiliary functions.  Its
converse and its unrestricted fibers have remained unclear.  We fix the two
diagonal values in the auxiliary annihilation recurrence as parameters and
work without a growth condition.  First, we disprove universal surjectivity.
At coupling $b=1/4$ and spin zero, Pillin's full pure-$A_4$ bootstrap family,
normalized by $\cK_0=\cK_2=0$ and $\cK_4=1$, has no auxiliary lift for any
fixed diagonal pair.  The obstruction is a four-particle calculation at two
intersecting annihilation divisors: all 25 generators of the homogeneous
recurrence module have zero regular transform value, and the lower-particle
kernel cannot repair it.  Second, for every fixed coupling, integer spin, and
fixed recurrence parameters, we classify the complete zero-transform kernel.
At each finite quotient stalk it is controlled by the explicit finite matrix
$C_x=B_xK_x$, where $K_x$ presents the transform syzygies and $B_x$ is the
reduced-divisor recurrence quotient.  Cartan's theorem B globalizes every
admitted step, and the full kernel is the direct sum of its even and odd
inverse-limit path spaces.  Every nonempty target fiber is an affine translate
of this kernel.  The nonsurjectivity theorem is specific to $b=1/4$; no
generic-coupling obstruction is asserted.
\end{abstract}

\subjclass[2020]{81T40, 32C35, 32E10, 37K40}
\keywords{Sinh--Gordon model, form factors, K-transform, coherent analytic sheaf, Cartan theorem B, inverse limit, nonsurjectivity}

\maketitle

\section{Introduction}
\label{sec:introduction}

The form-factor bootstrap turns factorized scattering data in
$1+1$-dimensional integrable quantum field theory into matrix elements of
local fields.  Its analytic foundations go back to the generalized form-factor
axioms of Karowski and Weisz, which rest in turn on the factorized
$S$-matrices of Zamolodchikov and Zamolodchikov, and to their subsequent
systematic development
\cite{Zamolodchikov_1979fs,Karowski_1978ef,Berg_1979co,Smirnov_1992}.  For the Sinh--Gordon model, explicit
solutions for fundamental fields, scalar operator towers, and general
symmetric polynomial families were obtained in
\cite{Fring_1993ff,Koubek_1993ot,Pillin_1998sg}.  These form factors now enter
rigorous work on correlation functions and the reconstruction of quantum-field
theoretic properties \cite{Kozlowski_2025om,Kozlowski_2025wa}.

A particularly useful construction writes the nonminimal scalar factor as a
finite sum over binary labels.  Entire, separately periodic functions
$p_n(\boldsymbol\beta\mid\boldsymbol\ell)$ satisfying a simple annihilation
recurrence are sent by this K-transform to symmetric meromorphic functions
$\cK_n(\boldsymbol\beta)$ satisfying the bootstrap residue equation.
Babujian--Karowski-type constructions establish the forward implication
\cite{Babujian_1999ef,Babujian_2002ef}; Kozlowski gives the precise Sinh--Gordon formulation
used below \cite[Prop.~2.3]{Kozlowski_2023ba}.  The paragraph following that
proposition leaves open whether all bootstrap solutions arise in this way.
The 2025 correlation-function formulation again states only the forward
implication, and its hypotheses add a growth bound on $p_n$ that is not
imposed here
\cite[Prop.~2.1]{Kozlowski_2025om}.  Free-field methods give one-to-one
correspondences in more restrictive generic descendant classes
\cite{Lukyanov_1995ff,Feigin_2009ff}, but they do not exhaust the unrestricted entire periodic
category considered here.

This paper gives two complementary answers.  The first is negative: the
universal surjectivity assertion is false.  At the self-dual value
$b=1/4$, a pure polynomial component in Pillin's all-level solution gives a
spin-zero target with
\begin{equation}
 \cK_0=\cK_2=0,
 \qquad \cK_4=1.
 \label{eq:intro-target}
\end{equation}
At the four-particle point
$(e^{\beta_1},\ldots,e^{\beta_4})=(-1,1,-i,i)$, every transform of a
homogeneous recurrence correction has zero regular value.  Changes of the
two-particle lift that remain recurrence-compatible also contribute zero.
Thus \eqref{eq:intro-target}, and hence the full Pillin family extending it,
has no lift.  A rank calculation shows that this conclusion holds for every
fixed choice of the two diagonal recurrence values, not just for the constant
choice usually adopted in examples.

The second answer describes all fibers once nonemptiness is known.  For
integer spin, periodicity and boost covariance descend the problem to the
finite quotient
\begin{equation}
 Z_n=\bigl((\C^*)^n/\C^*\bigr)/S_n.
 \label{eq:intro-quotient}
\end{equation}
Clearing the transform row gives a finite local coefficient map $A_x$.
Its syzygies form the exact fixed-particle kernel sheaf.  Quotienting boundary
values by the four-dimensional space of functions of the annihilating pair
gives a second finite map $B_x$.  If $K_x$ is any finite presentation of
$\ker A_x$, then
\begin{equation}
 C_x=B_xK_x
 \label{eq:intro-C}
\end{equation}
tests whether a lower kernel section has a recurrent zero-transform lift.
The admissible boundary image is $\im C_x$, and the homogeneous freedom is
$K_x(\ker C_x)$.  The quotient \eqref{eq:intro-quotient} is Stein, so Cartan's
theorem B removes the remaining global one-step obstruction.  Taking all
compatible steps separately in the two particle-number parities gives the
complete kernel; every nonempty fiber is its affine translate.

Four qualifications are structural.  The two diagonal recurrence values are
fixed parameters: changing them cannot generally be absorbed into a remainder
that is independent of all tail labels.  The inverse-limit theorem is an
integer-spin statement; noninteger spin gives zero source and target spaces by
periodicity.  The recurrence is imposed on the reduced annihilation divisor,
as in the original point-value condition, not on normal jets.  Finally, the
explicit empty fiber is proved at $b=1/4$.  An exact correction at the same
four-particle locus is nonzero for generic coupling, so the proof does not
support a generic-$b$ extrapolation.

Section~\ref{sec:setup} fixes the normalized function spaces and states the
main results.  Section~\ref{sec:fixed} constructs the fixed-particle transform
sheaves.  Section~\ref{sec:fibers} proves the inverse-limit fiber theorem.
Section~\ref{sec:target} builds the all-level pure-$A_4$ target, and
Section~\ref{sec:obstruction} proves that its fiber is empty.
Section~\ref{sec:scope} records low-particle checks and the generic-coupling
boundary.

\section{The normalized transform and main results}
\label{sec:setup}

Fix
\begin{equation}
 0<b<\frac12,
 \qquad \widehat b=\frac12-b,
 \qquad a=\sin(2\pi b).
 \label{eq:parameters}
\end{equation}
The scattering function is
\begin{equation}
 S(\beta)=
 \frac{\tanh(\beta/2-i\pi b)}{\tanh(\beta/2+i\pi b)}.
 \label{eq:S}
\end{equation}
We use the meromorphic continuation of the minimal factor
\begin{equation}
 F(\beta)=\exp\!\left\{-4\int_0^\infty
 \frac{\sinh(bx)\sinh(\widehat b x)\sinh(x/2)}
 {x\sinh^2x}
 \cos\!\left(\frac{x}{\pi}(i\pi-\beta)\right)dx\right\},
 \label{eq:F}
\end{equation}
initially defined for $0<\operatorname{Im}\beta<2\pi$.

\begin{definition}[Scalar target]
\label{def:target}
For $s\in\C$, an admissible scalar family is a sequence
$\cK=(\cK_n)_{n\ge0}$ such that $\cK_n$ is symmetric, separately
$2\pi i$-periodic and meromorphic, has only possible simple poles at
$\beta_i-\beta_j\in i\pi(1+2\Z)$, and satisfies
\begin{align}
 &\Res_{\beta_1=\beta_2+i\pi}\cK_n(\boldsymbol\beta_n)\,d\beta_1
 \notag\\
 &\quad=\frac{i}{F(i\pi)}
 \frac{1-\prod_{r=3}^nS(\beta_2-\beta_r)}
 {\prod_{r=3}^nF(\beta_2-\beta_r+i\pi)F(\beta_2-\beta_r)}
 \cK_{n-2}(\beta_3,\ldots,\beta_n).
 \label{eq:target-residue}
\end{align}
It has spin $s$ if
\begin{equation}
 \cK_n(\boldsymbol\beta_n+\theta\one_n)
 =e^{s\theta}\cK_n(\boldsymbol\beta_n)
 \qquad(\theta\in\C).
 \label{eq:target-spin}
\end{equation}
No growth condition is imposed.
\end{definition}

Let $g:\{0,1\}^2\to\C$ be symmetric and fixed, with
\begin{equation}
 g_{01}=g_{10}=\kappa:=-\frac1{aF(i\pi)},
 \qquad g_{00},g_{11}\in\C.
 \label{eq:g}
\end{equation}
The last two values are parameters throughout.

\begin{definition}[Labelled source]
\label{def:source}
An admissible labelled family is a sequence
\begin{equation}
 p_n:\C^n\times\{0,1\}^n\longrightarrow\C
 \end{equation}
whose rapidity dependence is entire and separately $2\pi i$-periodic,
which is invariant under simultaneous permutations of the pairs
$(\beta_j,\ell_j)$, and which obeys the spin covariance
\begin{equation}
 p_n(\boldsymbol\beta_n+\theta\one_n\mid\boldsymbol\ell_n)
 =e^{s\theta}p_n(\boldsymbol\beta_n\mid\boldsymbol\ell_n).
 \label{eq:p-spin}
\end{equation}
On $\beta_1=\beta_2+i\pi$ it satisfies
\begin{align}
 &p_n(\beta_2+i\pi,\beta_2,\beta_3,\ldots,\beta_n
       \mid\ell_1,\ell_2,\ell_3,\ldots,\ell_n)
 \notag\\
 &\quad=g(\ell_1,\ell_2)
 p_{n-2}(\beta_3,\ldots,\beta_n\mid\ell_3,\ldots,\ell_n)
 +h_n(\ell_1,\ell_2\mid\beta_2,\ldots,\beta_n),
 \label{eq:p-recurrence}
\end{align}
where $h_n$ is independent of every tail label
$\ell_3,\ldots,\ell_n$.  Conditions for all other pairs follow by
simultaneous permutation symmetry.
\end{definition}

Changing $g_{00}$ or $g_{11}$ changes Definition~\ref{def:source}.  Indeed,
the difference between two putative choices contributes
$(g'-g)(\ell_1,\ell_2)p_{n-2}$ to \eqref{eq:p-recurrence}; for a general
lower section this depends on the tail labels and cannot be absorbed into
$h_n$.  We therefore make no diagonal-gauge identification.

For $|\boldsymbol\ell|=\sum_j\ell_j$, define
\begin{align}
 T_n[p_n](\boldsymbol\beta_n)
  =\sum_{\boldsymbol\ell\in\{0,1\}^n}
  (-1)^{|\boldsymbol\ell|}
  \prod_{1\le i<j\le n}
  \left(1-i\frac{(\ell_i-\ell_j)a}
  {\sinh(\beta_i-\beta_j)}\right)
  p_n(\boldsymbol\beta_n\mid\boldsymbol\ell).
 \label{eq:transform}
\end{align}
The forward theorem states that $T[p]=(T_n[p_n])_{n\ge0}$ is an admissible
scalar family whenever $p$ is admissible
\cite[Prop.~2.3]{Kozlowski_2023ba}.

\begin{proposition}[Noninteger spin]
\label{prop:noninteger}
If $s\notin\Z$, both admissible spaces are zero.  Thus the transform is
trivially bijective in this branch.
\end{proposition}

\begin{proof}
Shift all $n\ge1$ rapidities by $2\pi i$.  Separate periodicity leaves either
a scalar or a labelled function unchanged, while covariance multiplies it by
$e^{2\pi i s}\ne1$.  Hence every positive-particle function vanishes.  At
$n=0$, covariance for every $\theta\in\C$ forces the constant to vanish unless
$s=0$.  The recurrence does not create a nonzero level from zero data.
\end{proof}

The first main result disproves the universal converse.

\begin{theorem}[An empty fiber]
\label{thm:empty}
Let $b=1/4$ and $s=0$.  For every fixed pair
$(g_{00},g_{11})\in\C^2$ with the off-diagonal value \eqref{eq:g}, there is a
full admissible scalar family $\cK^{A_4}$ such that
\begin{equation}
 \cK^{A_4}_0=\cK^{A_4}_2=0,
 \qquad \cK^{A_4}_4=1,
 \label{eq:A4-prefix}
\end{equation}
and
\begin{equation}
 T^{-1}(\cK^{A_4})=\varnothing.
\end{equation}
In particular, the unrestricted fixed-parameter K-transform is not universally
surjective.
\end{theorem}

The second result is valid at every fixed coupling.  Its notation is
constructed in Sections~\ref{sec:fixed} and \ref{sec:fibers}.

\begin{theorem}[Complete fixed-parameter fibers]
\label{thm:fibers-intro}
Fix $0<b<1/2$, $s\in\Z$, and one $g$ as in \eqref{eq:g}.  At each quotient
stalk $x$, the transform syzygies and the reduced-divisor recurrence have a
finite presentation
\begin{equation}
 A_x,\qquad B_x,\qquad C_x=B_xK_x,
 \label{eq:ABC-intro}
\end{equation}
where $\im K_x=\ker A_x$.  The complete zero-transform family space is the
direct sum of the even and odd inverse-limit path spaces defined by these
presentations.  If a scalar target $\cK$ has one admissible lift $p^*$, then
\begin{equation}
 T^{-1}(\cK)=p^*+\ker T.
 \label{eq:affine-intro}
\end{equation}
This describes all nonempty fibers; Theorem~\ref{thm:empty} supplies an empty
one.
\end{theorem}

\section{The fixed-particle coherent transform}
\label{sec:fixed}

We now assume $s\in\Z$.  Put $z_j=e^{\beta_j}$ and
$X_n=(\C^*)^n$.  The two reduced divisor unions relevant to
\eqref{eq:transform} are
\begin{equation}
 A_n^- =\bigcup_{i<j}\{z_i+z_j=0\},
 \qquad
 A_n^+ =\bigcup_{i<j}\{z_i-z_j=0\}.
 \label{eq:divisors}
\end{equation}
The first is the allowed annihilation divisor.  The second is an apparent
equal-rapidity pole of the individual coefficients, because
\begin{equation}
 \frac1{\sinh(\beta_i-\beta_j)}
 =\frac{2z_iz_j}{(z_i-z_j)(z_i+z_j)}.
 \label{eq:sinh-z}
\end{equation}

For a label word $\boldsymbol\ell$, abbreviate the coefficient in
\eqref{eq:transform} by
\begin{equation}
 c_{\boldsymbol\ell}(\boldsymbol\beta)
 =(-1)^{|\boldsymbol\ell|}
 \prod_{i<j}\left(1-i\frac{(\ell_i-\ell_j)a}
 {\sinh(\beta_i-\beta_j)}\right).
 \label{eq:row}
\end{equation}

\begin{lemma}[Equal-rapidity cancellation]
\label{lem:equal-cancel}
For every simultaneously symmetric labelled section $p$, the sum
$\sum_{\boldsymbol\ell}c_{\boldsymbol\ell}p_{\boldsymbol\ell}$ has no pole
on $A_n^+$.  The cancellation holds as a meromorphic identity on each whole
component, including its intersections with other divisors.
\end{lemma}

\begin{proof}
On $z_i=z_j$, pair a label word with the word obtained by interchanging
$\ell_i$ and $\ell_j$.  Their total signs agree.  The residue of the
$(i,j)$ factor changes sign, while all other factors are exchanged in pairs.
Simultaneous symmetry identifies the two boundary values of $p$.  Words with
$\ell_i=\ell_j$ have no pole.  Thus paired residues cancel as meromorphic
functions on the full component, not merely at its generic points.
\end{proof}

Common scaling of $\boldsymbol z$ is the boost action.  Let
\begin{equation}
 Y_n=X_n/\C^*\simeq(\C^*)^{n-1},
 \qquad Z_n=Y_n/S_n.
 \label{eq:YZ}
\end{equation}
The spin character defines a rank-$2^n$ labelled bundle
$E_{n,s}\to Y_n$ and a scalar line $L_{n,s}\to Y_n$.  The torus $Y_n$ is
affine and Stein.  Its finite permutation quotient exists as a normal analytic
space \cite{SHC_1953-1954__6__A12_0}; equivalently, its invariant Laurent
coordinate ring is finitely generated and $Z_n$ is affine.  Hence $Z_n$ is
Stein.

Write $\pi:Y_n\to Z_n$.  Taking the full permutation action, including the
spin and divisor linearizations, define the coherent sheaves
\begin{equation}
 \mathscr D_{n,s}=(\pi_*\mathscr E_{n,s})^{S_n},
 \qquad
 \mathscr M_{n,s}
 =\bigl(\pi_*(\mathscr L_{n,s}(A_n^-))\bigr)^{S_n}.
 \label{eq:DM}
\end{equation}
Their global sections are, respectively, the fixed-level labelled functions
in Definition~\ref{def:source} and symmetric spin-$s$ meromorphic functions
with at most simple annihilation poles.  Finite-group invariants are exact in
characteristic zero.  Lemma~\ref{lem:equal-cancel} therefore makes
\eqref{eq:transform} a coherent morphism
\begin{equation}
 \cT_{n,s}:\mathscr D_{n,s}\longrightarrow\mathscr M_{n,s}.
 \label{eq:sheaf-T}
\end{equation}
Set
\begin{equation}
 \mathscr I_{n,s}=\im\cT_{n,s},
 \qquad
 \mathscr N_{n,s}=\ker\cT_{n,s}.
 \label{eq:IN}
\end{equation}

The following local description is what makes \eqref{eq:IN} explicit.
Choose an ordered point $x\in Y_n$.  Let
\begin{equation}
 R=\cO_{Y_n,x},
 \qquad H_x\subseteq S_n,
 \qquad R_0=R^{H_x}\simeq\cO_{Z_n,\pi(x)},
 \label{eq:local-rings}
\end{equation}
where $H_x$ is the finite stabilizer.  Let $d_-$ and $d_+$ be products of
reduced local equations for the components of $A_n^-$ and $A_n^+$ through
$x$.  Each
\begin{equation}
 b_{\boldsymbol\ell}=d_-d_+c_{\boldsymbol\ell}
 \label{eq:b-row}
\end{equation}
is holomorphic.  Lemma~\ref{lem:equal-cancel} says that, for an invariant
label vector $p$, $\sum b_{\boldsymbol\ell}p_{\boldsymbol\ell}$ is divisible
by $d_+$.  In a local target frame define
\begin{equation}
 A_x(p)=
 \frac{\sum_{\boldsymbol\ell}b_{\boldsymbol\ell}
 p_{\boldsymbol\ell}}{d_+}.
 \label{eq:Ax}
\end{equation}
This is a finite $R_0$-linear map between the correctly linearized invariant
modules.

\begin{proposition}[Fixed-level image and kernel]
\label{prop:fixed-level}
At every quotient stalk,
\begin{equation}
 (\mathscr I_{n,s})_{\pi(x)}=\im A_x,
 \qquad
 (\mathscr N_{n,s})_{\pi(x)}=\ker A_x.
 \label{eq:fixed-stalks}
\end{equation}
Moreover,
\begin{equation}
 \im T_n=\Gamma(Z_n,\mathscr I_{n,s}),
 \qquad
 \ker T_n=\Gamma(Z_n,\mathscr N_{n,s}).
 \label{eq:fixed-global}
\end{equation}
On the free ordered locus away from $A_n^+$, one has $H_x=\{1\}$ and
$d_+=1$.  Thus $A_x$ is the row $(d_-c_{\boldsymbol\ell})$, which gives the
ordinary ideal and syzygy presentation
\begin{align}
 \mathscr I_{n,s}\otimes R
  &=\bigl(d_-c_{\boldsymbol\ell}:\boldsymbol\ell\in\{0,1\}^n\bigr),
 \label{eq:ideal-free}\\
 \mathscr N_{n,s}\otimes R
  &=\Syz_R\bigl(d_-c_{\boldsymbol\ell}\bigr).
 \label{eq:syzygy-free}
\end{align}
\end{proposition}

\begin{proof}
Equations \eqref{eq:b-row}--\eqref{eq:Ax} clear exactly the local target
frame $1/d_-$ and remove the apparent equal-rapidity factor.  Thus they give
the stalk map of \eqref{eq:sheaf-T}, including every stabilizer character,
and \eqref{eq:fixed-stalks} follows.

The exact sequence
\begin{equation}
 0\longrightarrow\mathscr N_{n,s}
 \longrightarrow\mathscr D_{n,s}
 \longrightarrow\mathscr I_{n,s}\longrightarrow0
 \label{eq:fixed-exact}
\end{equation}
consists of coherent sheaves on the Stein space $Z_n$.  Cartan's theorem B
gives $H^1(Z_n,\mathscr N_{n,s})=0$
\cite{SHC_1951-1952__4__A19_0}.  Taking global sections proves
\eqref{eq:fixed-global}.  When $H_x$ is trivial and $d_+=1$, the same map is
the row $(d_-c_{\boldsymbol\ell})$, giving
\eqref{eq:ideal-free}--\eqref{eq:syzygy-free}.
\end{proof}

Proposition~\ref{prop:fixed-level} does not assert that the image is the whole
allowed-pole line.  Indeed, the four-particle point used later is a base point
of every cleared generator.  Retaining the image ideal is essential when the
recurrence is imposed.

\section{The recurrence quotient and complete fibers}
\label{sec:fibers}

For an unordered pair $e=\{i,j\}$, let $D_e\subset Y_n$ be the reduced
annihilation component $z_i+z_j=0$.  In the label space
\begin{equation}
 V_n=\C^{\{0,1\}^n},
\end{equation}
let
\begin{equation}
 W_e=\{f(\boldsymbol\ell):f\text{ depends only on }
 (\ell_i,\ell_j)\}.
 \label{eq:We}
\end{equation}
This is four-dimensional and is precisely the allowed remainder space in
\eqref{eq:p-recurrence}.  If $\iota_e:D_e\hookrightarrow Y_n$ is the closed
inclusion, form the full edge sum upstairs and only then take invariants:
\begin{equation}
 \widetilde{\mathscr R}_{n,s}
 =\bigoplus_e(\iota_e)_*
 \bigl((\mathscr E_{n,s}|_{D_e})/\mathscr W_e\bigr),
 \qquad
 \mathscr R_{n,s}
 =\bigl(\pi_*\widetilde{\mathscr R}_{n,s}\bigr)^{S_n}.
 \label{eq:Rsheaf}
\end{equation}
This order is important when a stabilizer permutes annihilation components.
Restriction followed by the quotient defines
\begin{equation}
 \rho_{n,s}:\mathscr D_{n,s}\longrightarrow\mathscr R_{n,s}.
 \label{eq:rho}
\end{equation}

Let $q_{n-2}\in\Gamma(Z_{n-2},\mathscr N_{n-2,s})$.  On $D_e$, delete the
two coordinates in $e$, pull back $q_{n-2}$, extend it constantly in the two
deleted labels, multiply by $g(\ell_i,\ell_j)$, and take its class modulo
$W_e$.  These classes assemble to the invariant global section
\begin{equation}
 G_{n,s}(q_{n-2})
 =\bigl([g(\ell_i,\ell_j)\,\delta_e^*q_{n-2}]\bigr)_e
 \in\Gamma(Z_n,\mathscr R_{n,s}).
 \label{eq:G}
\end{equation}
For $n=2$ in a kernel family, we use the explicit convention
$G_{2,s}(q_0)=0$, since $q_0=0$.

\begin{lemma}[Recurrence as a quotient equality]
\label{lem:recurrence-quotient}
A pair $(q_{n-2},q_n)$ obeys the homogeneous form of
\eqref{eq:p-recurrence} for some tail-label-independent remainders if and only
if
\begin{equation}
 \rho_{n,s}(q_n)=G_{n,s}(q_{n-2}).
 \label{eq:rhoG}
\end{equation}
This equality is on the full reduced divisor union and therefore includes
shared, disjoint, and higher intersections.
\end{lemma}

\begin{proof}
On $D_e$, equality of the two quotient classes says exactly
\begin{equation}
 q_n-g(\ell_i,\ell_j)\delta_e^*q_{n-2}\in W_e,
\end{equation}
which is \eqref{eq:p-recurrence} with the lower and upper families replaced by
their difference.  Conversely, every allowed remainder lies in $W_e$ and
vanishes in the quotient.  Since \eqref{eq:Rsheaf} includes the whole reduced
component sheaf, equality retains all intersection germs.  It imposes values,
not normal jets or a nonreduced thickening.
\end{proof}

Restrict \eqref{eq:rho} to the transform syzygies and set
\begin{equation}
 \mathscr P_{n,s}=\ker(\rho_{n,s}|_{\mathscr N_{n,s}}),
 \qquad
 \mathscr J_{n,s}=\im(\rho_{n,s}|_{\mathscr N_{n,s}}).
 \label{eq:PJ}
\end{equation}
Thus
\begin{equation}
 0\longrightarrow\mathscr P_{n,s}
 \longrightarrow\mathscr N_{n,s}
 \xrightarrow{\rho_{n,s}}\mathscr J_{n,s}
 \longrightarrow0
 \label{eq:P-exact}
\end{equation}
is an exact sequence of coherent sheaves.

We next spell out the promised finite matrix.  At the ordered point $x$ of
\eqref{eq:local-rings}, let $E_x$ be the set of active annihilation edges and
choose a reduced local equation $d_e$ for each.  The recurrence block is
\begin{equation}
 B_x:p\longmapsto
 \left(p\bmod(d_e,W_e)\right)_{e\in E_x}
 \in\left[
 \bigoplus_{e\in E_x}(R/(d_e))\otimes(V_n/W_e)
 \right]^{H_x}.
 \label{eq:Bx}
\end{equation}
The direct sum is assembled before the $H_x$-invariants are taken.  Choose a
finite $R_0$-module presentation
\begin{equation}
 K_x:R_0^{m_x}\longrightarrow(R\otimes V_n)^{H_x},
 \qquad \im K_x=\ker A_x,
 \label{eq:Kx}
\end{equation}
and define
\begin{equation}
 C_x=B_xK_x.
 \label{eq:Cx}
\end{equation}
In a Walsh basis, $B_x$ simply selects the modes not supported inside $e$ and
reduces their coefficients modulo $d_e$.  Hence all three maps in
\eqref{eq:Cx} are finite analytic matrices determined directly by
\eqref{eq:row} and \eqref{eq:We}.

\begin{proposition}[Finite stalk criterion]
\label{prop:finite-stalk}
At every quotient stalk,
\begin{equation}
 (\mathscr N_{n,s})_{\pi(x)}=\ker A_x,
 \qquad
 (\mathscr J_{n,s})_{\pi(x)}=\im C_x,
 \qquad
 (\mathscr P_{n,s})_{\pi(x)}=K_x(\ker C_x).
 \label{eq:NJP-stalk}
\end{equation}
For prescribed boundary data $\eta_x$, a zero-transform recurrent lift exists
locally exactly when the finite system
\begin{equation}
 C_xu=\eta_x
 \label{eq:finite-system}
\end{equation}
is solvable, in which case $p_x=K_xu$ is one solution.
\end{proposition}

\begin{proof}
The first identity is Proposition~\ref{prop:fixed-level}.  Applying $B_x$ to
$\im K_x=\ker A_x$ gives the second.  For the third, if $p=K_xu$ and
$B_xp=0$, then $C_xu=0$; the converse is immediate.  Injectivity of $K_x$ is
not required.  Equation \eqref{eq:finite-system} is the same pair of
conditions written as
\begin{equation}
 \binom{A_x}{B_x}p=\binom0{\eta_x}.
\end{equation}
\end{proof}

\begin{theorem}[One-step lifting]
\label{thm:one-step}
For a fixed lower zero-transform section
$q_{n-2}\in\Gamma(Z_{n-2},\mathscr N_{n-2,s})$, the following are equivalent.
\begin{enumerate}[label=\textup{(\roman*)}]
 \item There is a global zero-transform $q_n$ satisfying the recurrence with
 $q_{n-2}$.
 \item $G_{n,s}(q_{n-2})\in\Gamma(Z_n,\mathscr J_{n,s})$.
 \item At every quotient stalk, \eqref{eq:finite-system} is solvable with
 $\eta_x=G_{n,s}(q_{n-2})_x$.
\end{enumerate}
When these conditions hold and $q_n^*$ is one lift, every lift is
\begin{equation}
 q_n^*+k_n,
 \qquad k_n\in\Gamma(Z_n,\mathscr P_{n,s}).
 \label{eq:one-step-torsor}
\end{equation}
\end{theorem}

\begin{proof}
By Proposition~\ref{prop:fixed-level}, a zero-transform section is a global
section of $\mathscr N_{n,s}$.  Lemma~\ref{lem:recurrence-quotient} makes its
recurrence exactly the prescribed image under $\rho_{n,s}$.  This proves the
necessity of (ii), and Proposition~\ref{prop:finite-stalk} makes (ii)
equivalent to the stalkwise condition (iii).

For sufficiency, apply global sections to \eqref{eq:P-exact}.  The space
$Z_n$ is Stein and $\mathscr P_{n,s}$ is coherent, so Cartan B gives
$H^1(Z_n,\mathscr P_{n,s})=0$.  Therefore every global section of
$\mathscr J_{n,s}$ lifts to one of $\mathscr N_{n,s}$.  The kernel of the
global map is $\Gamma(Z_n,\mathscr P_{n,s})$, proving
\eqref{eq:one-step-torsor}.
\end{proof}

The full answer is the set of infinite paths through these one-step
correspondences.  Define
\begin{equation}
 \cC_{n,s}=\left\{(u,v)\in
 \Gamma(\mathscr N_{n-2,s})\times\Gamma(\mathscr N_{n,s}):
 \rho_{n,s}(v)=G_{n,s}(u)\right\}
 \label{eq:correspondence}
\end{equation}
and, for $\epsilon\in\{0,1\}$,
\begin{align}
 \fK_s^{(\epsilon)}
 =\bigl\{(q_{\epsilon+2m})_{m\ge0}:{}&
 q_{\epsilon+2m}\in\Gamma(\mathscr N_{\epsilon+2m,s}),
 \notag\\[-2mm]
 & (q_{\epsilon+2m-2},q_{\epsilon+2m})
 \in\cC_{\epsilon+2m,s}\text{ for }m\ge1\bigr\}.
 \label{eq:path-space}
\end{align}
At level zero the transform is the identity whenever the scalar can be
nonzero, so every kernel path has $q_0=0$.

\begin{theorem}[Complete kernel]
\label{thm:complete-kernel}
For every fixed $b$, integer $s$, and fixed $g$ as in \eqref{eq:g},
\begin{equation}
 \ker T=\fK_s^{(0)}\oplus\fK_s^{(1)},
 \qquad q_0=0.
 \label{eq:full-kernel}
\end{equation}
Equivalently, at each step retain a branch precisely when
$G_{n,s}(q_{n-2})$ lies in the coherent image $\mathscr J_{n,s}$; if it does,
all next terms are one chosen lift plus
$\Gamma(Z_n,\mathscr P_{n,s})$.  Only infinite retained paths are admissible.
\end{theorem}

\begin{proof}
An admissible zero-transform family restricts at each level to a section of
$\mathscr N_{n,s}$, and Lemma~\ref{lem:recurrence-quotient} places each
successive pair in \eqref{eq:correspondence}.  Conversely, an infinite path
consists of entire periodic symmetric spin sections, has zero transform at
every level, and supplies an allowed remainder on every annihilation
component.  Particle number changes by two, so the two parities are
independent.  The recursive description follows from
Theorem~\ref{thm:one-step}.
\end{proof}

\begin{corollary}[All nonempty fibers]
\label{cor:fiber}
If $p^*$ is one admissible lift of a target $\cK$ in the same fixed recurrence
category, then
\begin{equation}
 T^{-1}(\cK)=p^*+
 \bigl(\fK_s^{(0)}\oplus\fK_s^{(1)}\bigr).
 \label{eq:fiber}
\end{equation}
\end{corollary}

\begin{proof}
The difference of two lifts has zero transform, and the difference of their
remainders remains tail-label independent.  It therefore belongs to
\eqref{eq:full-kernel}.  Conversely, adding any kernel path preserves every
linear admissibility condition and the transform.
\end{proof}

This is an exact inverse-limit parametrization, not a finite-dimensional
basis.  Without a growth restriction, the global section spaces appearing in
\eqref{eq:path-space} already contain arbitrary holomorphic functions on
rapidity tori.  The finite content is local: every one-step membership and
homogeneous freedom is given by \eqref{eq:Cx}.

\section{An all-level pure-\texorpdfstring{$A_4$}{A4} target}
\label{sec:target}

We next construct the target in Theorem~\ref{thm:empty}.  Pillin uses an
effective coupling $B$ and
\begin{equation}
 S_{\mathrm P}(\theta)
 =\frac{i\sinh\theta+\sin(\pi B/2)}
 {i\sinh\theta-\sin(\pi B/2)}.
 \label{eq:Pillin-S}
\end{equation}
The elementary identity
\begin{equation}
 \frac{\tanh(\theta/2-i\alpha)}
 {\tanh(\theta/2+i\alpha)}
 =\frac{\sinh\theta-i\sin(2\alpha)}
 {\sinh\theta+i\sin(2\alpha)}
 \label{eq:tanh-identity}
\end{equation}
shows that \eqref{eq:Pillin-S} equals \eqref{eq:S} under
\begin{equation}
 B=4b.
 \label{eq:B-map}
\end{equation}
After the substitution $x=2t$ in the minimal-factor integral, the two
normalizations are related by
\begin{equation}
 F^{\min}_{\mathrm P}(\theta)=\frac{F(\theta)}{F(i\pi)}.
 \label{eq:F-map}
\end{equation}

Pillin's polynomial ansatz is
\begin{equation}
 \mathcal F_n^{\mathrm P}(\boldsymbol\theta)
 =H_nQ_n(\boldsymbol x)
 \prod_{i<j}\frac{F^{\min}_{\mathrm P}(\theta_i-\theta_j)}
 {x_i+x_j},
 \qquad x_j=e^{\theta_j},
 \label{eq:Pillin-ansatz}
\end{equation}
Thus the scalar factor left after removing the minimal two-particle factors is
\begin{equation}
 \mathcal F_n^{\mathrm P}
 =\left(\prod_{i<j}F(\theta_i-\theta_j)\right)\cK_n,
 \qquad
 \cK_n=\frac{H_n}{F(i\pi)^{n(n-1)/2}}
 \frac{Q_n(\boldsymbol x)}{\prod_{i<j}(x_i+x_j)}.
 \label{eq:Pillin-to-K}
\end{equation}
and its kinematic residue equation is equivalent to
\begin{equation}
 Q_{n+2}(-x,x,\boldsymbol x_n)
 =(-1)^nD_n(x\mid\boldsymbol x_n)Q_n(\boldsymbol x_n),
 \label{eq:Pillin-recursion}
\end{equation}
where, with $\omega=e^{i\pi B/2}$,
\begin{align}
 D_n(x\mid\boldsymbol x_n)
 =\frac{x}{2(\omega-\omega^{-1})}
 \bigg[&\prod_{j=1}^n(x+\omega x_j)(x-\omega^{-1}x_j)
 \notag\\
 &-\prod_{j=1}^n(x-\omega x_j)(x+\omega^{-1}x_j)\bigg].
 \label{eq:Dn}
\end{align}
Pillin proves an all-level symmetric solution of
\eqref{eq:Pillin-recursion}; its new component at level $r$ is
\begin{equation}
 Q_{r,r}(x_1,\ldots,x_r)=\prod_{1\le i<j\le r}(x_i+x_j)
 \label{eq:Pillin-kernel}
\end{equation}
\cite[Eqs.~(15), (19)]{Pillin_1998sg}.

\begin{proposition}[Pure-$A_4$ family]
\label{prop:A4}
At $b=1/4$, there is a full spin-zero admissible target
$\cK^{A_4}$ satisfying \eqref{eq:A4-prefix}, with all odd levels zero.
\end{proposition}

\begin{proof}
Under \eqref{eq:B-map}, $b=1/4$ gives $B=1$ and $\omega=i$.  In Pillin's
independent-component expansion, set every free coefficient except $A_4$ to
zero.  Then
\begin{equation}
 Q_{2,4}=0,
 \qquad
 Q_{4,4}=\prod_{1\le i<j\le4}(x_i+x_j),
 \label{eq:Q4}
\end{equation}
and the all-level recursion constructs $Q_{2m,4}$ for every $m\ge2$.
Define
\begin{equation}
 \cK^{A_4}_{2m}
 =C_{2m}\frac{Q_{2m,4}}
 {\prod_{i<j}(x_i+x_j)},
 \qquad
 \cK^{A_4}_{2m+1}=0,
 \label{eq:A4-family}
\end{equation}
where $C_{2m}$ contains $H_{2m}$ and the powers of $F(i\pi)$ from
\eqref{eq:F-map}.  Pillin's residue normalization fixes their ratios, and one
overall nonzero rescaling makes $\cK_4^{A_4}=1$.

The numerator degree is $n(n-1)/2$, equal to the pair-sum denominator degree,
so \eqref{eq:A4-family} has spin zero.  It is symmetric and separately
periodic, and its only possible poles are simple pair-sum poles.  Equations
\eqref{eq:Pillin-recursion}--\eqref{eq:Dn} were derived from the same
kinematic residue axiom after the minimal-factor separation
\eqref{eq:Pillin-ansatz}; the convention matches by
\eqref{eq:B-map}--\eqref{eq:F-map}.  Thus \eqref{eq:target-residue} holds at
every even level.  Equation \eqref{eq:Q4} gives
$\cK_0^{A_4}=\cK_2^{A_4}=0$ and $\cK_4^{A_4}=1$.
\end{proof}

The all-level input in Proposition~\ref{prop:A4} is the proved polynomial
recursion, not an extrapolation from a finite prefix.  We now show that its
four-particle value cannot come from the labelled transform.

\section{The four-particle obstruction}
\label{sec:obstruction}

Throughout this section $b=1/4$, so $a=1$.  Work at the ordered point
\begin{equation}
 z^*=(-1,1,-i,i).
 \label{eq:zstar}
\end{equation}
Only the annihilation components $D_{12}$ and $D_{34}$ are active.  Use
transverse coordinates
\begin{equation}
 z_1=-(1+\xi),\quad z_2=1,
 \qquad
 z_3=-i(1+\eta),\quad z_4=i,
 \label{eq:transverse}
\end{equation}
and write $d_{12}=z_1+z_2=-\xi$ and
$d_{34}=z_3+z_4=-i\eta$.  Split a label word as
$A=(\ell_1,\ell_2)$ and $C=(\ell_3,\ell_4)$.

Let $R$ be the ordered local ring at \eqref{eq:zstar}.  The homogeneous
recurrence condition for an upper correction $q_4$ is
\begin{equation}
 q_4\bmod d_{12}\in W_{12},
 \qquad
 q_4\bmod d_{34}\in W_{34}.
 \label{eq:two-boundaries}
\end{equation}

\begin{lemma}[The 25-generator module]
\label{lem:25module}
The local module defined by \eqref{eq:two-boundaries} is
\begin{equation}
 \mathscr H_{4,0,z^*}
 =R\one+d_{12}(R\otimes W_{34})
 +d_{34}(R\otimes W_{12})
 +d_{12}d_{34}(R\otimes\C^{16}).
 \label{eq:25module}
\end{equation}
At \eqref{eq:zstar}, the regular value of the original transform row on every
displayed generator is zero.
\end{lemma}

\begin{proof}
The two four-dimensional spaces consist of functions of disjoint label pairs,
so
\begin{equation}
 W_{12}\cap W_{34}=\C\one.
 \label{eq:W-intersection}
\end{equation}
Subtract the common intersection value from the two boundary terms and divide
the remaining $D_{12}$ and $D_{34}$ contributions by their reduced local
equations.  The residual section vanishes on both components and is divisible
by $d_{12}d_{34}$.  This gives \eqref{eq:25module}.

Let $e_A$ and $e_C$ denote the characteristic functions of a head or tail
label pair.  Substitution of \eqref{eq:transverse} into \eqref{eq:sinh-z} and
the coefficient row \eqref{eq:row} gives the complete table.
\begin{equation}
\begin{array}{c@{\qquad}c@{\qquad}c}
\toprule
\text{generator family}&\text{multiplicity}&
 \left.T_4[\text{generator}]\right|_{z^*}^{\mathrm{reg}}\\
\midrule
\one&1&0\\
d_{12}e_C&C\in\{0,1\}^2\ (4)&0\\
d_{34}e_A&A\in\{0,1\}^2\ (4)&0\\
d_{12}d_{34}e_Ae_C&(A,C)\in\{0,1\}^4\ (16)&0\\
\bottomrule
\end{array}
\label{eq:25table}
\end{equation}
The forward residue identity with zero lower transform shows that each
displayed transform is holomorphic along both active divisors.  Consequently,
$R$-linearity reduces the regular-value calculation for
\eqref{eq:25module} to the 25 listed generators.
For completeness, the last line has an immediate factor check.  A label word
that does not cut both active edges retains a factor $d_{12}$ or $d_{34}$ and
vanishes.  For each of the four proper cuts, the cross-edge product contains,
up to a nonzero sign,
\begin{equation}
 1+\frac{a^2}{\sinh^2(i\pi/2)}=1+\frac1{i^2}=0.
 \label{eq:mixed-zero}
\end{equation}
For a one-boundary generator, only the two cuts of the cleared active edge
survive; their signs and cross factors are opposite at \eqref{eq:zstar}.
The constant generator is their sum over the remaining pair.  This proves all
four rows of \eqref{eq:25table}.  There are $1+4+4+16=25$ generators.
\end{proof}

Lemma~\ref{lem:25module} says that changing $p_4$ while holding a lower lift
fixed cannot change the regular value at $z^*$.  We must also exclude changing
the lower $p_2$ itself.  At $b=1/4$, write
\begin{equation}
 u=\frac{i}{\sinh(\beta_1-\beta_2)}.
\end{equation}
The two-particle transform row in label order $(00,01,10,11)$ is
\begin{equation}
 (1,-1-u,-1+u,1).
 \label{eq:T2-row}
\end{equation}

\begin{lemma}[Symmetry-closed lower block]
\label{lem:lower-block}
Let $q=(q_{00},q_{01},q_{10},q_{11})^T$ be a local two-particle section with
$T_2[q]=0$ that can occur below a recurrent four-particle section.  The
transform equation, the shared-pair recurrence equations, and their
simultaneous coordinate-label swap give
\begin{equation}
 M_g(u)q=0,
 \label{eq:Mgq}
\end{equation}
where
\begin{equation}
 M_g(u)=
 \begin{pmatrix}
 1&-1-u&-1+u&1\\
 \kappa&-\kappa&-g_{00}&g_{00}\\
 g_{11}&-g_{11}&-\kappa&\kappa\\
 \kappa&-g_{00}&-\kappa&g_{00}\\
 g_{11}&-\kappa&-g_{11}&\kappa
 \end{pmatrix}.
 \label{eq:Mg}
\end{equation}
If $(g_{00},g_{11})\ne(\kappa,\kappa)$, its kernel is the
label-constant line.  If $g_{00}=g_{11}=\kappa$, the kernel is generated by
\begin{equation}
 v^{(0)}=(1,\tfrac12,\tfrac12,0)^T,
 \qquad
 v^{(1)}=(0,\tfrac12,\tfrac12,1)^T.
 \label{eq:v01}
\end{equation}
\end{lemma}

\begin{proof}
The first row is \eqref{eq:T2-row}.  Comparing restrictions to two
annihilation components that share one coordinate eliminates the
tail-independent remainders and gives the next two rows; applying the
simultaneous swap gives the last two.  Direct $3\times3$ minors of
\eqref{eq:Mg} include
\begin{equation}
 \pm(g_{00}-\kappa)^2,
 \qquad
 \pm(g_{11}-\kappa)^2.
 \label{eq:square-minors}
\end{equation}
The vector $(1,1,1,1)^T$ is always in the kernel.  If either difference in
\eqref{eq:square-minors} is nonzero, the rank is at least three and hence is
exactly three.  This includes the nonconstant locus
$g_{00}g_{11}=\kappa^2$, where the unsymmetrized three-row block alone would
have rank two.  At the constant pair the matrix has rank two, and direct
substitution shows that \eqref{eq:v01} is a basis of its kernel.
\end{proof}

\begin{proposition}[No four-level lift]
\label{prop:no-prefix-lift}
For every fixed diagonal pair, a recurrent spin-zero prefix with
$T_0[p_0]=T_2[p_2]=0$ satisfies
\begin{equation}
 T_4[p_4](z^*)=0.
 \label{eq:T4-zero}
\end{equation}
Consequently it cannot transform to \eqref{eq:A4-prefix}.
\end{proposition}

\begin{proof}
Since $T_0$ is the identity at spin zero, $p_0=0$.  First suppose the diagonal
pair is nonconstant.  Lemma~\ref{lem:lower-block} makes $p_2$ label
independent, say $p_2=f$.  It has the exact recurrent upper lift $p_4=0$:
on component $e$ choose
\begin{equation}
 h_{4,e}(\ell_i,\ell_j)=-g(\ell_i,\ell_j)f.
 \label{eq:zero-lift}
\end{equation}
This is independent of the tail labels.  Every other lift differs by a section
of \eqref{eq:25module}, whose transform value is zero by
Lemma~\ref{lem:25module}.  Hence \eqref{eq:T4-zero} holds.

Now let $g_{00}=g_{11}=\kappa$.  A recurrent extension of a lower kernel
section, modulo \eqref{eq:25module}, is
\begin{equation}
 q_{A,C}(\xi,\eta)=\kappa\bigl(v_A(\xi)+v_C(\eta)\bigr).
 \label{eq:lower-extension}
\end{equation}
For $v=\alpha v^{(0)}+\delta v^{(1)}$, direct contraction with the original
row gives
\begin{equation}
 \kappa R(\xi,\eta)
 \bigl[(\alpha-\delta)(\xi)-(\alpha-\delta)(\eta)\bigr],
 \label{eq:R-contraction}
\end{equation}
where
\begin{equation}
 R(\xi,\eta)=
 \frac{4(1+\xi)(1+\eta)(2+\xi+\eta)
 (\xi\eta+i\xi-i\eta)}
 {(\xi^2+2\xi+2)(\eta^2+2\eta+2)
 (\xi^2+2\xi+\eta^2+2\eta+2)}.
 \label{eq:R-function}
\end{equation}
Thus the regular value is zero.  The apparent extra two-particle syzygy
\begin{equation}
 v^{(\sinh)}=
 \left(0,\frac{i-\sinh\beta_{12}}2,
 \frac{i+\sinh\beta_{12}}2,0\right)^T
 \label{eq:vsinh}
\end{equation}
would give the nonzero value $-2i\kappa$, but its shared-pair defect is
$-i$; it does not satisfy \eqref{eq:Mgq}.  Hence every admissible lower change
is covered by \eqref{eq:v01}--\eqref{eq:R-function}, and
\eqref{eq:T4-zero} again follows.
\end{proof}

\begin{proof}[Proof of Theorem~\ref{thm:empty}]
Proposition~\ref{prop:A4} supplies a full admissible target whose
four-particle function is the constant one.  Any hypothetical full lift would
in particular give a recurrent prefix through level four.  Proposition
\ref{prop:no-prefix-lift} evaluates its transform at $z^*$ as zero, whereas
$\cK_4^{A_4}(z^*)=1$.  This contradiction is independent of the fixed
diagonal pair.
\end{proof}

\section{Low-particle controls and scope}
\label{sec:scope}

The abstract sheaves reproduce the smallest levels directly.  At $n=0$ the
kernel is zero.  At $n=1$ the row is $(1,-1)$, so
\begin{equation}
 q_1(\beta)=c e^{s\beta}(1,1)
 \label{eq:n1-kernel}
\end{equation}
is the fixed-level kernel.  At $n=2$, with
$u=ia/\sinh(\beta_{12})$, the row is
\begin{equation}
 (1,-1-u,-1+u,1).
 \label{eq:T2-general}
\end{equation}
Away from its divisors it has a rank-three syzygy module.  The recurrence
quotient is zero at $n=2$ because every label is an endpoint label.  At
$n=3$, the quotient has label rank four, and the prescription induced by
\eqref{eq:n1-kernel} is tail independent.  At the first disjoint
four-particle stratum,
\begin{equation}
 \dim(W_{12}\cap W_{34})=1,
\end{equation}
which is exactly the intersection used in Lemma~\ref{lem:25module}.

Independent exact-arithmetic controls were used to check the finite algebra
that supports, but does not replace, the all-level proofs.
\begin{table}[ht]
\caption{Definition-level verification controls.  The all-level ingredients
are Proposition~\ref{prop:A4} and Theorems~\ref{thm:one-step} and
\ref{thm:complete-kernel}; the finite rows below check their first nontrivial
instances.}
\label{tab:checks}
\centering
\small
\begin{tabularx}{\textwidth}{>{\raggedright\arraybackslash}p{0.28\textwidth}X>{\raggedright\arraybackslash}p{0.18\textwidth}}
\toprule
Object & Exact check & Outcome\\
\midrule
Pure-$A_4$ target & Componentwise polynomial recursion, permutation symmetry,
and homogeneity at $n=4,6,8$ & All checks pass\\
Recurrence-zero stalk & Laurent expansion of all 25 generators in
\eqref{eq:25module}, including pole cancellation & All regular values vanish\\
Lower-lift repair & Three exact $T_2$ syzygies and their induced $T_4$ values;
shared-pair defects & Values $0,0,-2i$; defects $0,0,-i$\\
Arbitrary diagonal pair & Symbolic minors of \eqref{eq:Mg}, the
determinant-zero nonconstant locus, and 32 disjoint recurrence equations &
Ranks $3,3,2$ in the three regimes\\
Complete kernel & Original rows at $n=0,1,2$, recurrence quotients through
$n=4$, and the edge-swapping stabilizer & Finite presentations agree with
\eqref{eq:NJP-stalk}\\
\bottomrule
\end{tabularx}
\end{table}

The obstruction should not be extrapolated away from $b=1/4$.  To make this
boundary concrete, write
\begin{equation}
 t=e^{2\pi ib}=c+ia,
 \qquad c=\cos(2\pi b),
\end{equation}
and use the analogous point $(-1,1,-t,t)$.  For a label word $\ell$ and an
unordered pair $e=\{i,j\}$, let $e^c=\{k,l\}$ with $k<l$, and define
\begin{equation}
 Q_\ell(z)=\sum_{i<j}(z_i+z_j)
 \prod_{r\in e,\,s\in e^c}(z_r+z_s)
 (\ell_k-\ell_l)(z_k-z_l).
 \label{eq:generic-Q}
\end{equation}
On $z_i+z_j=0$, only the term indexed by $e^c$ survives, and that term depends
only on $(\ell_i,\ell_j)$.  Hence $Q$ is a symmetric homogeneous recurrence-zero
section.  Multiplication by
\begin{equation}
 H(z)=\left(\sum_{r=1}^4z_r^{-2}\right)^3
\end{equation}
makes it spin zero.  Direct evaluation gives
\begin{equation}
 T_4[HQ](-1,1,-t,t)=4096\,i a^3c^4.
 \label{eq:generic-value}
\end{equation}
This is nonzero whenever $c\ne0$.  Thus the constant-germ cokernel used in
Section~\ref{sec:obstruction} disappears at this same locus for generic
coupling.  Equation \eqref{eq:generic-value} neither constructs an all-level
lift nor proves generic surjectivity; it only rules out reusing the
$b=1/4$ point-value proof unchanged.

There is a second scope boundary.  Cartan B supplies holomorphic lifts but no
growth estimates.  Consequently Theorem~\ref{thm:complete-kernel} answers the
unrestricted conditions in Definition~\ref{def:source}; it does not by itself
classify the polynomial-growth subclass used for correlation estimates in
\cite{Kozlowski_2025om}.  Likewise, the quotient \eqref{eq:Rsheaf} encodes the
reduced-divisor value recurrence.  A jet-level or nonreduced recurrence would
require a different quotient sheaf.

\section{Conclusion}
\label{sec:conclusion}

The Sinh--Gordon K-transform is not a universal parametrization of the
unrestricted bootstrap families.  At $b=1/4$, the all-level pure-$A_4$
target has an empty fiber for every fixed diagonal recurrence pair.  The
failure is visible at four particles, but it is not a truncated-target
artifact: Pillin's recursion supplies the full admissible family.

For every fixed coupling and fixed recurrence parameters, all integer-spin
fibers nevertheless have a complete common form.  The original transform row
and the tail-independence quotient give finite stalk matrices
$A_x$, $B_x$, and $C_x=B_xK_x$.  Cartan B turns stalkwise image membership
into a global one-step lift, and the two parity inverse limits give the exact
kernel.  Every target fiber is therefore either empty or an affine translate
of this explicitly presented path space.

The remaining image problem is parameter dependent.  In particular,
\eqref{eq:generic-value} shows that the exceptional four-particle value
obstruction does not persist at the same locus for generic $b$.  Determining
the image there requires a different stratum, higher normal jets, or an
all-level compatibility obstruction.  Adding a growth condition would pose a
separate analytic lifting problem because the Stein-space argument used here
does not control growth.

\section*{Data and code availability}

No empirical data were used, and no step of any proof above is computer
assisted.  The all-level results rest on
Proposition~\ref{prop:A4} and Theorems~\ref{thm:one-step} and
\ref{thm:complete-kernel}, proved analytically here; the finite algebra
recorded in Table~\ref{tab:checks} is stated and derived in the text, at
\eqref{eq:25module}, \eqref{eq:Mg} and \eqref{eq:NJP-stalk}, and each entry
can be recomputed by hand from the displayed coefficient row \eqref{eq:row}.
Exact rational and symbolic arithmetic was run independently while preparing
the paper, and it reproduced those entries; because it certifies a
transcription rather than an argument, nothing above depends on it.

\bibliographystyle{amsplain}
\bibliography{references,unverified}

\end{document}
